\hypertarget{objetivos}{%
\section{Objetivos}\label{objetivos}}

\begin{frame}{Objetivos}
\protect\hypertarget{objetivos-1}{}
\begin{itemize}
\item
  Apresentar o formato textual WebAssembly (WAT) e o modelo de execu\c{c}\~{a}o
  baseado em máquina de pilha.
\item
  Definir as regras de tradu\c{c}\~{a}o de express\~{o}es e comandos da IRT para
  instru\c{c}\~{o}es WAT.
\item
  Discutir a reconstru\c{c}\~{a}o do controle de fluxo estruturado de
  WebAssembly a partir dos saltos explícitos da IRT.
\end{itemize}
\end{frame}

\hypertarget{introduuxe7uxe3o}{%
\section{Introdu\c{c}\~{a}o}\label{introduuxe7uxe3o}}

\begin{frame}[fragile]{WebAssembly como Alvo}
\protect\hypertarget{webassembly-como-alvo}{}
\[\text{IRT} \xrightarrow{\text{opt}} \text{IRT otimizada} \xrightarrow{\text{codegen}} \text{Módulo WAT} \xrightarrow{\texttt{wat2wasm}} \text{Binário Wasm}\]

\begin{itemize}
\tightlist
\item
  \textbf{WebAssembly} (Wasm): especifica\c{c}\~{a}o de máquina virtual com
  sem\^{a}ntica formal (W3C)
\item
  Excelente para fins pedagógicos:

  \begin{itemize}
  \tightlist
  \item
    Sem\^{a}ntica \textbf{precisa} e publicada
  \item
    Representa\c{c}\~{a}o textual (WAT) \textbf{legível}
  \item
    Ferramentas: \passthrough{\lstinline!wabt!} (compilar) e
    \passthrough{\lstinline!wasmtime!} (executar)
  \end{itemize}
\end{itemize}
\end{frame}

\begin{frame}[fragile]{Por que WebAssembly?}
\protect\hypertarget{por-que-webassembly}{}
\begin{itemize}
\tightlist
\item
  Executado no browser (Chrome, Firefox, Safari) e fora dele (wasmtime,
  wasmer)
\item
  Modelo de seguran\c{c}a: sandbox + verifica\c{c}\~{a}o de tipos estática
\item
  Tipo único no nosso subconjunto: \passthrough{\lstinline!i32!}
  (inteiro de 32 bits)
\item
  Todos os valores da IRT --- inteiros, ponteiros, endere\c{c}os --- mapeiam
  para \passthrough{\lstinline!i32!}
\end{itemize}
\end{frame}

\hypertarget{o-formato-wat}{%
\section{O Formato WAT}\label{o-formato-wat}}

\begin{frame}[fragile]{Estrutura de um Módulo WAT}
\protect\hypertarget{estrutura-de-um-muxf3dulo-wat}{}
\begin{lstlisting}
(module
  (import "env" "print"    (func $print    (param i32)))
  (import "env" "read_int" (func $read_int (result i32)))
  (memory 4)
  (global $heap_ptr (mut i32) (i32.const 0))
  (func $main
    ;; corpo da fun\c{c}\~{a}o
  )
  (export "main" (func $main))
)
\end{lstlisting}

\begin{itemize}
\tightlist
\item
  \passthrough{\lstinline!(memory 4)!} = 4 páginas de 64 KB = 256 KB de
  memória linear
\item
  \passthrough{\lstinline!$heap\_ptr!} = ponteiro de aloca\c{c}\~{a}o
  (\emph{bump-pointer})
\item
  Fun\c{c}\~{o}es \passthrough{\lstinline!print!} e
  \passthrough{\lstinline!read\_int!} importadas do ambiente
\end{itemize}
\end{frame}

\begin{frame}{Gramática do Subconjunto WAT}
\protect\hypertarget{gramuxe1tica-do-subconjunto-wat}{}
\[\begin{array}{lcl}
\mathit{mod} & ::= & \mathtt{(module}\ \vec{\mathit{import}}\ \mathit{mem}?\ \vec{\mathit{global}}\ \vec{\mathit{func}}\ \vec{\mathit{export}}\mathtt{)} \\
\mathit{func} & ::= & \mathtt{(func}\ \$f\ \vec{\mathit{param}}\ \mathit{result}?\ \vec{\mathit{local}}\ \vec{i}\mathtt{)} \\
i & ::= & \mathtt{i32.const}\ n \mid \mathtt{local.get}\ \$x \mid \mathtt{local.set}\ \$x \\
  & \mid & \mathit{binop} \mid \mathtt{i32.load} \mid \mathtt{i32.store} \\
  & \mid & \mathtt{call}\ \$f \mid \mathtt{drop} \mid \mathtt{return} \\
  & \mid & \mathtt{(block}\ \$\ell\ \vec{i}\mathtt{)} \mid \mathtt{(loop}\ \$\ell\ \vec{i}\mathtt{)} \\
  & \mid & \mathtt{(if\ (then\ \vec{i})\ (else\ \vec{i}))} \\
  & \mid & \mathtt{br}\ \$\ell \mid \mathtt{br\_if}\ \$\ell
\end{array}\]
\end{frame}

\hypertarget{semuxe2ntica-operacional}{%
\section{Sem\^{a}ntica Operacional}\label{semuxe2ntica-operacional}}

\begin{frame}{Máquina de Pilha}
\protect\hypertarget{muxe1quina-de-pilha}{}
WebAssembly é uma \textbf{máquina de pilha}: instru\c{c}\~{o}es consomem do topo
e empilham resultados.

Estado de uma fun\c{c}\~{a}o: \((\sigma, \mu)\)

\begin{itemize}
\tightlist
\item
  \(\sigma = [v_1, \ldots, v_n]\) --- pilha de operandos (topo \`{a}
  direita)
\item
  \(\mu\) --- variáveis locais (incluindo par\^{a}metros)
\end{itemize}

Regras de avalia\c{c}\~{a}o ---
\((\sigma, \mu) \xrightarrow{i} (\sigma', \mu')\):

\[\frac{}{{(\sigma, \mu) \xrightarrow{\mathtt{i32.const}\ n} (\sigma \cdot n,\ \mu)}} \qquad \frac{\mu(\$x) = v}{(\sigma, \mu) \xrightarrow{\mathtt{local.get}\ \$x} (\sigma \cdot v,\ \mu)}\]

\[\frac{\sigma = \sigma' \cdot v_2 \cdot v_1 \quad v = \mathit{eval}(\mathit{op}, v_1, v_2)}{(\sigma, \mu) \xrightarrow{\mathit{binop}} (\sigma' \cdot v,\ \mu)}\]
\end{frame}

\begin{frame}[fragile]{Controle de Fluxo Estruturado}
\protect\hypertarget{controle-de-fluxo-estruturado}{}
WebAssembly usa controle de fluxo \textbf{estruturado} (sem
\passthrough{\lstinline!goto!}):

\begin{longtable}[]{@{}
  >{\raggedright\arraybackslash}p{(\columnwidth - 2\tabcolsep) * \real{0.5000}}
  >{\raggedright\arraybackslash}p{(\columnwidth - 2\tabcolsep) * \real{0.5000}}@{}}
\toprule\noalign{}
\begin{minipage}[b]{\linewidth}\raggedright
Construtor
\end{minipage} & \begin{minipage}[b]{\linewidth}\raggedright
Comportamento de \passthrough{\lstinline!br!}
\end{minipage} \\
\midrule\noalign{}
\endhead
\passthrough{\lstinline!(block $l ...)!} & Salta para o \textbf{fim} do
bloco (\emph{break}) \\
\passthrough{\lstinline!(loop $l ...)!} & Salta para o \textbf{início}
do la\c{c}o (\emph{continue}) \\
\passthrough{\lstinline!(if (then ...) (else ...))!} & Consome topo;
executa ramo verdadeiro se \(\neq 0\) \\
\passthrough{\lstinline!br\_if $l!} & Consome topo; ramifica se
\(\neq 0\) \\
\bottomrule\noalign{}
\end{longtable}
\end{frame}

\begin{frame}[fragile]{Padr\~{a}o Can\^{o}nico de While em WAT}
\protect\hypertarget{padruxe3o-canuxf4nico-de-while-em-wat}{}
\begin{lstlisting}
(block $exit
  (loop $head
    ;; avalia\c{c}\~{a}o da condi\c{c}\~{a}o de saída
    <condi\c{c}\~{a}o>
    i32.eqz          ;; nega\c{c}\~{a}o
    br_if $exit      ;; sai se condi\c{c}\~{a}o falsa
    ;; corpo do la\c{c}o
    <corpo>
    br $head         ;; volta ao início
  )
)
\end{lstlisting}
\end{frame}

\hypertarget{traduuxe7uxe3o-da-irt-para-wat}{%
\section{Tradu\c{c}\~{a}o da IRT para
WAT}\label{traduuxe7uxe3o-da-irt-para-wat}}

\begin{frame}[fragile]{Alocador de Memória (Bump-Pointer)}
\protect\hypertarget{alocador-de-memuxf3ria-bump-pointer}{}
\begin{lstlisting}
(func $alloc (param $n i32) (result i32)
  (local $base i32)
  global.get $heap_ptr
  local.set $base
  global.get $heap_ptr
  local.get $n
  i32.const 8
  i32.mul
  i32.add
  global.set $heap_ptr
  local.get $base
)
\end{lstlisting}

\(\mathit{alloc}(n)\): aloca \(n\) palavras de 8 bytes, retorna endere\c{c}o
base.
\end{frame}

\begin{frame}{Tradu\c{c}\~{a}o de Express\~{o}es
\(\mathcal{E}\llbracket \cdot \rrbracket\)}
\protect\hypertarget{traduuxe7uxe3o-de-expressuxf5es-mathcalellbracket-cdot-rrbracket}{}
\[\mathcal{E}\llbracket \mathbf{CONST}(n) \rrbracket = [\mathtt{i32.const}\ n]\]

\[\mathcal{E}\llbracket \mathbf{TEMP}(t) \rrbracket = [\mathtt{local.get}\ \$t]\]

\[\mathcal{E}\llbracket \mathbf{BINOP}(\mathit{op}, e_1, e_2) \rrbracket = \mathcal{E}\llbracket e_1 \rrbracket \cdot \mathcal{E}\llbracket e_2 \rrbracket \cdot [\mathit{op}_\mathtt{i32}]\]

\[\mathcal{E}\llbracket \mathbf{MEM}(e) \rrbracket = \mathcal{E}\llbracket e \rrbracket \cdot [\mathtt{i32.load}]\]

\[\mathcal{E}\llbracket \mathbf{CALL}(\mathbf{NAME}\ f, \vec{e}) \rrbracket = \mathcal{E}\llbracket \vec{e} \rrbracket \cdot [\mathtt{call}\ \$f]\]
\end{frame}

\begin{frame}[fragile]{Tabela de Operadores}
\protect\hypertarget{tabela-de-operadores}{}
\begin{longtable}[]{@{}llll@{}}
\toprule\noalign{}
IRT & WAT & IRT & WAT \\
\midrule\noalign{}
\endhead
ADD & \passthrough{\lstinline!i32.add!} & EQ &
\passthrough{\lstinline!i32.eq!} \\
SUB & \passthrough{\lstinline!i32.sub!} & NEQ &
\passthrough{\lstinline!i32.ne!} \\
MUL & \passthrough{\lstinline!i32.mul!} & LT &
\passthrough{\lstinline!i32.lt\_s!} \\
DIV & \passthrough{\lstinline!i32.div\_s!} & LE &
\passthrough{\lstinline!i32.le\_s!} \\
MOD & \passthrough{\lstinline!i32.rem\_s!} & GT &
\passthrough{\lstinline!i32.gt\_s!} \\
AND & \passthrough{\lstinline!i32.and!} & GE &
\passthrough{\lstinline!i32.ge\_s!} \\
OR & \passthrough{\lstinline!i32.or!} & & \\
\bottomrule\noalign{}
\end{longtable}
\end{frame}

\begin{frame}{Tradu\c{c}\~{a}o de Comandos
\(\mathcal{S}\llbracket \cdot \rrbracket\)}
\protect\hypertarget{traduuxe7uxe3o-de-comandos-mathcalsllbracket-cdot-rrbracket}{}
\[\mathcal{S}\llbracket \mathbf{MOVE}(\mathbf{TEMP}\ t, e) \rrbracket = \mathcal{E}\llbracket e \rrbracket \cdot [\mathtt{local.set}\ \$t]\]

\[\mathcal{S}\llbracket \mathbf{MOVE}(\mathbf{MEM}(a), e) \rrbracket = \mathcal{E}\llbracket a \rrbracket \cdot \mathcal{E}\llbracket e \rrbracket \cdot [\mathtt{i32.store}]\]

\[\mathcal{S}\llbracket \mathbf{RETURN}([e]) \rrbracket = \mathcal{E}\llbracket e \rrbracket \cdot [\mathtt{return}]\]
\end{frame}

\begin{frame}{Reconstru\c{c}\~{a}o de Controle de Fluxo}
\protect\hypertarget{reconstruuxe7uxe3o-de-controle-de-fluxo}{}
A IRT usa JUMP/CJUMP/LABEL; WAT usa block/loop/if.

\textbf{Reconhecimento de while}: padr\~{a}o na sequ\^{e}ncia linearizada:
\[\mathbf{LABEL}\ \ell_h \;\; \mathbf{CJUMP}(e, \ell_b, \ell_f) \;\; \mathbf{LABEL}\ \ell_b \;\; \text{corpo} \;\; \mathbf{JUMP}(\ell_h) \;\; \mathbf{LABEL}\ \ell_f\]

\textbf{Reconhecimento de if/else}: padr\~{a}o:
\[\mathbf{CJUMP}(e, \ell_t, \ell_f) \;\; \mathbf{LABEL}\ \ell_t \;\; \text{then} \;\; \mathbf{JUMP}(\ell_e) \;\; \mathbf{LABEL}\ \ell_f \;\; \text{else} \;\; \mathbf{LABEL}\ \ell_e\]
\end{frame}

\hypertarget{conclusuxe3o}{%
\section{Conclus\~{a}o}\label{conclusuxe3o}}

\begin{frame}[fragile]{Sumário do Capítulo}
\protect\hypertarget{sumuxe1rio-do-capuxedtulo}{}
\begin{itemize}
\tightlist
\item
  WebAssembly é uma máquina de \textbf{pilha} com tipo único
  \passthrough{\lstinline!i32!}
\item
  Controle de fluxo \textbf{estruturado}: block, loop, if (sem goto)
\item
  Pipeline completo: \textbf{TImp → IRT → WAT → Wasm}
\item
  Tradu\c{c}\~{a}o de express\~{o}es: empilhar operandos, depois opera\c{c}\~{a}o
\item
  Tradu\c{c}\~{a}o de la\c{c}os: \emph{linearizar} IRT, \emph{reconhecer} padr\~{a}o,
  \emph{reconstruir} block/loop
\end{itemize}
\end{frame}

\begin{frame}[fragile]{Próximos Passos}
\protect\hypertarget{pruxf3ximos-passos}{}
\begin{itemize}
\tightlist
\item
  \textbf{Tipos múltiplos}: Wasm 2.0 adiciona
  \passthrough{\lstinline!f32!}, \passthrough{\lstinline!f64!},
  \passthrough{\lstinline!v128!} (SIMD)
\item
  \textbf{Aloca\c{c}\~{a}o de memória avan\c{c}ada}: GC, ponteiros managed
\item
  \textbf{Otimiza\c{c}\~{o}es específicas de Wasm}: tail calls, instru\c{c}\~{o}es SIMD
\item
  \textbf{Pipeline completo}: integrar todos os capítulos em um
  compilador funcional
\end{itemize}
\end{frame}
